% u03-beamer-fontspec: a beamer deck with fontspec's fonts (TeX Gyre Heros)
% and unicode-math, overlays, blocks and columns
\documentclass{beamer}
\usetheme{Madrid}
\usepackage{fontspec}
\usepackage{unicode-math}
\setsansfont{texgyreheros-regular.otf}[
  BoldFont=texgyreheros-bold.otf,
  ItalicFont=texgyreheros-italic.otf]
\setmathfont{latinmodern-math.otf}
\title{A Unicode deck}
\author{FlashTeX}
\date{11 October 2026}
\begin{document}
\begin{frame}
  \titlepage
\end{frame}
\begin{frame}{Overlays}
  \begin{itemize}
    \item<1-> First point, with “quotes” and an em dash — here.
    \item<2-> Second point: $∑_{k=1}^{n} k = \frac{n(n+1)}{2}$.
    \item<3-> Third point, \alert{alerted}.
  \end{itemize}
\end{frame}
\begin{frame}{Blocks and columns}
  \begin{columns}
    \column{0.5\textwidth}
    \begin{block}{Definition}
      A set $A ⊆ ℝ$ is open when every point has a neighbourhood in $A$.
    \end{block}
    \column{0.5\textwidth}
    \begin{exampleblock}{Example}
      The interval $(0, 1)$ is open; $[0, 1]$ is not.
    \end{exampleblock}
  \end{columns}
\end{frame}
\begin{frame}{Math}
  \[ ∫_{-∞}^{∞} e^{-x^2}\,dx = \sqrt{π} \]
  Some text after the display, long enough to wrap onto a second line of
  the frame in the theme's sans serif font.
\end{frame}
\end{document}
